\section{Introduction}
Cooperative trajectory planning enables multiple unmanned aerial vehicles (UAVs) to execute coordinated missions in a shared three-dimensional workspace. When several vehicles approach an intersection or a narrow passage, their trajectories must satisfy obstacle clearance, mutual separation, and motion limits simultaneously. Aerial-swarm research has established the importance of these coupled requirements~\cite{chung2018}, while agile quadrotor teams and coordinated formation transitions demonstrate their practical relevance~\cite{kushleyev2012,preiss2017}. Recent work combines three-dimensional reciprocal velocity obstacles with multi-agent reinforcement learning to coordinate collision avoidance within structured urban airspace~\cite{zhong2025rvo}.

Computational efficiency is essential when a feasible coordinated motion must be obtained within a limited planning interval. Parallel stitching addresses this requirement in online parking replanning by evaluating candidate connecting trajectories in parallel and preserving continuity with the ongoing motion~\cite{li2022stitching}. Enhanced sequential convex programming has been integrated with model predictive control to address formation tracking and correct linearization errors that can compromise feasibility~\cite{liu2025escp}. Conflict-free multi-robot planning coupled with predictive control connects coordination decisions to feasible motion~\cite{wu2025conflict}. Cooperative exploration combines dynamic Voronoi partitioning with learned obstacle avoidance~\cite{zhao2025exploration}, while multi-UAV coverage planning jointly considers task allocation and travel time under repeated round-trip operation~\cite{li2025coverage}. Together, these studies motivate methods that account for coupled coordination, safety, and timing requirements while keeping computation manageable.

High-conflict regions arise where assigned routes intersect or several UAVs compete for a shared passage around obstacles. Safe corridors express collision avoidance through conservative linear restrictions. A useful reference can nevertheless produce local corridors that conflict with one another or with bounded motion. Under a prescribed completion time, all of these requirements must be reconciled within the original flight horizon. A failed local solve therefore calls for a mechanism that can revise the coordinated motion and its surrounding geometry while preserving the task boundaries.

This paper develops a lightweight iterative optimization method (LIOM) for fixed-horizon, rest-to-rest cooperative flight. The planner first attempts a hard relative-corridor QP. An accepted result is returned immediately. Otherwise, a relaxed formulation supplies intermediate candidates for corridor reconstruction. A second recovery layer handles candidates whose reconstructed motion violates a physical limit, tightening the kinematic tolerance and attempting exact finishing. The double-integrator model yields affine defects, so both stages remain convex QPs with linear constraints. The recovery loop treats kinematic defects and relative-corridor slack separately. Exact control integration and complete-interval checking connect the numerical candidates to the prescribed physical task.

The main contributions are twofold:
\begin{enumerate}
\item A fixed-horizon cooperative formulation combines reachable-speed corridor padding with an affine two-penalty QP. The deadline and rest-to-rest endpoints remain fixed throughout refinement. An explicit feasible-point construction establishes nonemptiness of the relaxed local problem, while conservative polyhedral motion bounds and box-distance screening retain sparse linear constraints.
\item An acceptance-driven solution procedure combines a hard-QP entry, capped continuation, corridor reconstruction, and endpoint moment correction. Failed continuous-motion checks trigger tighter recovery and an exact-kinematics finish in the current corridors. Complete-interval checks on the integrated output determine acceptance. This procedure concentrates additional recovery effort on attempts that require it.
\end{enumerate}
The evaluation follows the three computational paths: direct hard-QP solution, relaxed corridor recovery, and acceptance-driven finishing. Matched hard formulations receive the same reconstruction opportunities. Independently checked winding tasks expose local corridor conflicts, while scaling, parameter, and tracking studies characterize computation and execution.
